% Chapter 13 of Deep Learning as an Experimental Science: extrapolation beyond the scales measured.
\labchapter{Extrapolation}{Predicting the large from the small}{On these tasks, larger language models do worse.}{The received reading of the eleven winning tasks of the Inverse Scaling Prize (McKenzie et al.\ 2023)}{Narrowed: on a larger family of models only four of the eleven tasks still scaled inversely; a trend over a few decades of scale describes those decades.}{ch:extrapolation}
\begin{refsection}[chapters/13-extrapolation.bib]

\begin{chaptersummary}
In 2022 a public contest collected eleven tasks on which language models did worse the larger they were, across model families spanning five orders of magnitude of training compute. On a larger family with five times the compute, six of the curves turned upward, one rose throughout, and one solved example in the prompt changed the other four: the trends were real over the range measured and did not hold beyond it. The method treats a forecast at a new scale as a measurement with three sources of error: the uncertainty of the fitted parameters, which grows, far from the data, linearly with the distance; the error of the functional form, which grows quadratically and is nearly invisible inside the data; and a change of regime, which no fit can detect. It checks a forecasting procedure by backtesting it on held-out large scales. In the laboratory, a straight line fitted in log--log coordinates to the six smallest of thirteen widths underestimates the loss of the largest by a factor of 2.8, and lines fitted to five to eight rungs have 95\% intervals that cover 21\% of held-out runs beyond 1.2 decades, against 88\% for a law with a floor; on digits with noisy labels, a miniature inverse-scaling trend reverses within a decade.
\end{chaptersummary}

\section{The claim: eleven tasks on which bigger was worse}\label{sec:extrapolation:claim}

By 2022 the loss of a language model on held-out text was one of the most predictable numbers in deep learning. Fitted power laws described it over many orders of magnitude of parameters, data and compute \parencite{KaplanEtAl2020,HoffmannEtAl2022}, laboratories chose the size of their largest runs by fitting smaller ones (\cref{ch:fitting,ch:protocol}), and a lower loss usually came with better scores on downstream tasks. Large models thus combined a predictable loss with abilities that were harder to anticipate, the combination that \textcite{GanguliEtAl2022} called predictability and surprise.\sidenote{The surprise usually meant an ability that appeared abruptly with scale; \cref{ch:emergence} examines how much of the abruptness belongs to the metric.}

Exceptions were known. On TruthfulQA, a set of questions that some people answer falsely because of a common misconception, the largest models tested were generally the least truthful \parencite{LinHiltonEvans2022}. A task on which performance gets worse as models get larger, or as their loss on the training objective improves, is said to show \term{inverse scaling}\index{inverse scaling}.

In June 2022 a group of researchers, most of them at New York University and FAR AI, opened a public contest for such tasks, the Inverse Scaling Prize, which offered up to \$100,000 for ``an important task where larger language models do worse'' \parencite{McKenzieEtAl2022}. Participants submitted a dataset of at least 300 examples with plots over the sizes of GPT-3, and the organizers evaluated each submission on further model families, among them private ones from Anthropic and DeepMind, so that a task could not owe its trend to one family. Two rounds, with deadlines on 27 August and 27 October 2022, drew 99 submissions. The grand and second prizes went unawarded, mainly because no task convincingly demonstrated real-world consequences of the failure; eleven tasks received third prizes of \$5,000. The organizers and the task authors reported the results together \parencite{McKenzieEtAl2023}.

\begin{finding}[\textcite{McKenzieEtAl2023}]\label{fd:extrapolation:prize}
On the eleven datasets collected by the contest, task performance got worse as training compute grew, over all or part of the range, and for most of the tasks the trend held across most of the model series evaluated: GPT-2, GPT-3, OPT, pretrained Anthropic models, Gopher and Chinchilla, with instruction-tuned variants, from about $10^{18}$ to $10^{23}$ training FLOPs. The authors grouped the tasks by four suspected causes: a prior from pretraining that overrides the instruction in the prompt; imitation of undesirable patterns in the training data; an easy distractor task that the model performs instead of the real one; and few-shot demonstrations that are correct but misleading.
\end{finding}

The tasks are easy for people. In Memo Trap a model must write a phrase that begins like a famous quotation but ends differently; the larger the model, the more often it completes the quotation. In Hindsight Neglect the demonstrations show bets whose outcomes agree with their expected values, and the question is about a bet whose outcome does not; larger models increasingly judged the bet by its outcome.

The contest asked for tasks on which larger models do worse, and its winners were presented as such tasks. Read as a statement about scale rather than about the models measured, the result became the claim of this chapter. It mattered beyond the eleven tasks: if scale could make some behaviors worse, then scaling could also encourage behavior at odds with human preferences, the implication with which \textcite{WeiEtAl2023} open their study of the tasks.

The organizers built the safeguards that such a claim needs: held-out model families, labels checked by crowd workers, and results with few-shot prompts and across training checkpoints as well as zero-shot. Their final report, written after the first post-mortem, is more careful than the reception: it describes U-shaped and inverted-U trends and concludes that scaling trends are ``less reliable at predicting the behavior of larger-scale models than previously understood.'' The claim was reasonable when the tasks were collected. The trends held across families with different data and architectures, several had a plausible mechanism, such as a prior from pretraining that grows stronger with scale, and the field's experience with loss curves suggested that trends in scale are smooth and monotone. The expectation that made the inverse trends alarming also made them look durable.

\section{The experiment: a dose--response curve read beyond its last dose}\label{sec:extrapolation:experiment}

Each winning task was a small experiment with scale as its treatment.

\begin{experiment}[a winning task of the Inverse Scaling Prize, 2022]
\field{question}{Does performance on the task get worse as language models get larger?}
\field{units}{pretrained models, one per size, from each of several families}
\field{treatment}{scale, measured by training compute (or parameters) within a family}
\field{levels}{the sizes each family offered, e.g.\ Gopher 44M--280B parameters; about $10^{18}$--$10^{23}$ FLOPs in all}
\field{held fixed}{the task's prompts, answer options and scoring; zero-shot in the main analysis}
\field{varied}{across families: training data, architecture, tokenizer, ratio of data to parameters}
\field{response}{accuracy, taking the most probable answer option (loss for one task)}
\field{replicates}{one model per size; at least 300 items per task}
\field{analysis}{scaling plots per family, reviewed against a rubric: roughly monotone inverse scaling on public models and inverse scaling on some private ones}
\end{experiment}

This is a dose--response design: the dose is scale and the response is accuracy. It can show that over the range measured the response fell; the held-out families make a quirk of one family's data or code unlikely, and the human validation makes wrong labels unlikely.

It cannot show what happens at a larger dose. A family contributed a handful of models over a few decades of compute, one per size, with no estimate of how far a second model of the same size would have landed, and the trend was judged from the plots. ``Larger models do worse'' is a statement about models larger than the largest one measured, and the design measured none of them. Two further features would matter. The prompt and the scoring were held fixed, so each trend belonged to the task under one setting of the instrument (\cref{ch:instruments}). And the families differ in their ratio of data to parameters, so their curves cannot be laid end to end on one axis of compute, as the authors themselves caution.\sidenote{The axis is itself an estimate: for several families the organizers computed training FLOPs as $6ND$ from the parameters and tokens, and the counts for Gopher and Chinchilla were provided privately.}

\section{The post-mortem: the curves turn}\label{sec:extrapolation:postmortem}

\paragraph{A larger family.} In November 2022 Jason Wei, Najoung Kim, Yi Tay and Quoc Le, at Google, ran the four first-round tasks on PaLM, and later all eleven; Kim was also an author of the prize's report \parencite{WeiEtAl2023}. They used the 8B, 62B and 540B models of the original PaLM paper and a 1B model trained by the same protocol, with the contest's prompts and scoring. The largest used about $2.5\times10^{24}$ training FLOPs, five times as much as Chinchilla 70B, the model trained with the most compute in the contest.

\begin{finding}[\textcite{WeiEtAl2023}]\label{fd:extrapolation:ushape}
On PaLM models of 1B, 8B, 62B and 540B parameters, six of the eleven tasks scaled in a U shape, performance falling and then rising up to the largest model (Hindsight Neglect, Sig Figs, Resisting Correction, Modus Tollens, Memo Trap, NeQA); one improved throughout (Repetitive Algebra); four still scaled inversely (Pattern Match Suppression, Redefine, Into the Unknown, Prompt Injection). With one solved example added to each prompt, those four became U-shaped or flat, and chain-of-thought prompting changed four of the nine classification tasks that were U-shaped or flat under one-shot prompting to positive scaling, by the authors' count.
\end{finding}

Their explanation of this \term{U-shaped scaling}\index{U-shaped scaling} is the third cause of \cref{fd:extrapolation:prize}, the distractor task: small models can do neither the real task nor the distractor and answer near chance, medium models learn the distractor and do worse, and large models learn to set it aside. The authors call the hypothesis speculative.\sidenote{They also point to earlier U-shaped curves among the tasks of the BIG-bench collection, TruthfulQA among them.}

Four points limit what the post-mortem shows, and its authors state most of them. First, each curve has four points, one model per size, and in three of the six U shapes the upturn is the 540B model alone. Second, an upturn is not a recovery: on four of the six tasks the 540B model still scored below the 1B model, and with one-shot prompts it was behind or not substantially ahead of the 1B model on six of the ten tasks that were then U-shaped or flat. Third, PaLM is another family, trained on other data, and its curve is not the continuation of Gopher's or GPT-3's. Fourth, the measurement depends on the prompt. The first version of the paper had modified some of the prompts and reported three of the four first-round tasks as U-shaped; the second version, posted eleven days later after feedback from Ethan Perez, one of the prize's organizers, used the contest's prompts unchanged and found two.\sidenote{Both versions are on arXiv (2211.02011, versions 1 and 2); the comment on version 2 records the change of prompts. The final version, with all eleven tasks, appeared at EMNLP 2023.}

\paragraph{One more model.} The GPT-4 technical report added a point \parencite{OpenAI2023}: on Hindsight Neglect, accuracy fell across three sizes of GPT-3 and GPT-3.5, then reached 100\% with GPT-4, a result the report filed under capabilities that ``remain hard to predict.'' McKenzie et al., who received GPT-4 results on five further tasks, report improvement on three, mixed results on Pattern Match Suppression and poor ones on Memo Trap, and they are uncertain whether GPT-4 should count as a member of the same series as its predecessors.

\paragraph{Reversals both ways.} Trends also turned for the worse. On Prompt Injection, McKenzie et al.\ found performance that improved with scale and then declined, an inverted U. Whether a trend persists beyond the decades measured, and even its direction there, cannot be read from the curve inside them.

\paragraph{How common is an unreliable trend?} \Textcite{LourieHuCho2025} re-examined the 46 downstream tasks of \textcite{GadreEtAl2025}, whose models span 0.011 to 6.9 billion parameters, and classified how accuracy changed as the validation loss improved: 18 tasks, or 39\%, improved smoothly and predictably, and the others scaled inversely, non-monotonically, noisily, without a trend, or by an abrupt breakthrough. The classification is by inspection, as the authors say. Changing the corpus on which the validation loss is measured, or the experimental setup in which the same task is evaluated, could reverse which pretraining corpus looked better or turn a clean scaling law into a non-monotone one.

\Textcite{CaballeroEtAl2023} proposed a \term{smoothly broken power law}\index{broken power law}, a chain of power laws joined by smooth breaks. On the held-out large scales of their benchmark it extrapolated better than four other forms, with the lowest error on 69\% of the downstream vision tasks and 75\% of the language tasks, and it can bend into a double-descent curve or a sharp inflection. It cannot anticipate a break that has not happened: on four-digit addition they could extrapolate only from points at or past the break, and they conclude that a sufficiently sharp break well beyond the largest fitted point cannot at present be extrapolated.

\paragraph{The counterpoint: forecasts that held.} Some forecasts were made in advance and held. For GPT-4, a power law with an irreducible term, $L(C)=aC^b+c$, fitted to models trained by the same method with at most a ten-thousandth of the compute, predicted the final loss on an internal code corpus, shortly after the run started and without its partial results; a forecast of the pass rate on groups of HumanEval problems, from models with at most a thousandth of the compute, was registered before training finished and held except on the easiest group \parencite{OpenAI2023}. The Llama 3 team forecast the 405B model's accuracy on the ARC Challenge in two steps, from compute to the loss on the correct answers and from that loss to accuracy, across four orders of magnitude, and slightly underestimated it \parencite{LlamaTeam2024}. \Textcite{GadreEtAl2025} predicted the validation loss of two large runs to within 0.7\% from experiments with 300 times less compute, and their average top-1 error on 17 tasks to within 0.05\% and 3.6\% from 20 times less; individual tasks were much noisier, and without the 1.4B model in the fit one of the errors rose from 0.05\% to 10.64\%. Ladders with about 1\% of the target compute predicted the accuracy of 7B and 13B models to within two points on four multiple-choice tasks \parencite{BhagiaEtAl2025}, and fits to about a hundred public models predicted several abilities that had looked abrupt \parencite{RuanMaddisonHashimoto2024}.

The forecasts that held share three features: the quantity was smooth, a loss or an average over many tasks, or a single task reached through such a quantity; the ladder and the target were trained alike, or their differences were modeled; and the forecast was checked against the result. The ones that failed predicted one task's metric beyond the range measured, often across a change of family, prompt or regime.

\begin{verdict}{narrowed}
Over the range the contest measured, the inverse trends were real: on most of the eleven tasks, performance fell with scale across most of the families tested. On a larger family with five times the compute, six of the curves turned upward and one rose throughout, a single example in the prompt changed the other four, and GPT-4 solved one task outright. What survives is a description of the decades measured, for the families and prompts used. The step from ``falls across this range'' to ``larger models do worse'' was an extrapolation, and nothing in the design could support it.
\end{verdict}

\section{The method: uncertainty that grows with distance}\label{sec:extrapolation:variance}

A forecast at a new scale is a measurement at a point where nothing was measured. In a scaling study the \term{design} is the ladder of scales trained\index{ladder of scales}, $x_1,\dots,x_n$, where $x$ is the logarithm of the number of parameters, examples or FLOPs, and the response $y$ is the logarithm of a loss or of a transformed metric; a forecast predicts the response of one new run at a scale $x_0$. It is an \term{interpolation}\index{interpolation} when $x_0$ lies within the range of the design and an \term{extrapolation}\index{extrapolation} when it lies beyond.

\paragraph{Four sources of error.} A forecast misses for four reasons. The new run has its own noise, the spread of runs at one scale (\cref{ch:seeds}), which no forecast removes. The fitted parameters are uncertain, and their uncertainty, carried to $x_0$, grows with the distance from the data, linearly once the distance is large (\cref{thm:extrapolation:variance}). The functional form may be wrong, and its error, almost invisible inside the data, grows quadratically (\cref{prop:extrapolation:misspec}). And the process may change: a mechanism that does not act inside the range measured, such as a distractor that only large models learn to ignore, bends the curve beyond it (\cref{rem:extrapolation:break}). Prediction intervals account for the first two sources, a test of the form for part of the third, and nothing computed from the data for the fourth.

The first two are a theorem about the straight line, the form that a power law takes in logarithmic coordinates.

\begin{theorem}[Prediction at a new scale]\label{thm:extrapolation:variance}\index{prediction interval}
Let $x_1,\dots,x_n$ be design points, $n\ge2$, not all equal, and let $y_i=a+bx_i+\varepsilon_i$ with uncorrelated errors of mean zero and variance $\sigma^2$. Write $\bar x$ for the mean of the design points, $S_{xx}=\sum_i(x_i-\bar x)^2$, and $\hat y(x_0)=\hat a+\hat bx_0$ for the least-squares prediction at a point $x_0$.
\begin{enumerate}
\item The prediction is unbiased for $a+bx_0$, the slope has $\Var\hat b=\sigma^2/S_{xx}$, and
\begin{equation}\label{eq:extrapolation:var}
\Var\hat y(x_0)=\sigma^2\Bigl(\frac1n+\frac{(x_0-\bar x)^2}{S_{xx}}\Bigr).
\end{equation}
For a new observation $y_0=a+bx_0+\varepsilon_0$ whose error has mean zero and variance $\sigma^2$ and is uncorrelated with $\varepsilon_1,\dots,\varepsilon_n$, the error $y_0-\hat y(x_0)$ has mean zero and variance $\sigma^2\bigl(1+1/n+(x_0-\bar x)^2/S_{xx}\bigr)$.
\item If every $x_i$ lies in an interval $[L,U]$, then $S_{xx}\le n(U-L)^2/4$, so that $\Var\hat b\ge4\sigma^2/\bigl(n(U-L)^2\bigr)$, with equality if and only if $n$ is even and half of the $x_i$ equal $L$ and half equal $U$.
\item If moreover $x_0=L+t(U-L)$ with $t>1$, then $\Var\hat y(x_0)\ge(2t-1)^2\sigma^2/n$, with equality if and only if every $x_i$ equals $L$ or $U$ and a fraction $t/(2t-1)$ of them equal $U$.
\end{enumerate}
\end{theorem}

\begin{proof}
The plan is to write the prediction as the true line plus a linear combination of the errors, which gives part~(i), and then to bound $S_{xx}$ for designs inside $[L,U]$ by one inequality, which gives parts~(ii) and~(iii).

\emph{Part (i).} Put $w_i=(x_i-\bar x)/S_{xx}$, so that $\sum_iw_i=0$, $\sum_iw_ix_i=1$ and $\sum_iw_i^2=1/S_{xx}$. The least-squares line passes through the point of means with slope $\hat b=\sum_iw_iy_i$, so $\hat y(x_0)=\bar y+\hat b(x_0-\bar x)$. Substituting $y_i=a+bx_i+\varepsilon_i$ and using the two sums of the $w_i$, we get $\hat b=b+\sum_iw_i\varepsilon_i$ and $\bar y=a+b\bar x+\bar\varepsilon$, hence
\[
\hat y(x_0)=a+bx_0+\bar\varepsilon+(x_0-\bar x)\sum_iw_i\varepsilon_i ,
\]
whose mean is $a+bx_0$. Because the errors are uncorrelated with variance $\sigma^2$, the sum $\sum_iw_i\varepsilon_i$ has variance $\sigma^2\sum_iw_i^2=\sigma^2/S_{xx}$, which is $\Var\hat b$, the mean $\bar\varepsilon$ has variance $\sigma^2/n$, and their covariance is $(\sigma^2/n)\sum_iw_i=0$. Hence $\Var\hat y(x_0)=\sigma^2/n+(x_0-\bar x)^2\sigma^2/S_{xx}$, which is \eqref{eq:extrapolation:var}. The new error $\varepsilon_0$ is uncorrelated with $\hat y(x_0)-a-bx_0$, a combination of $\varepsilon_1,\dots,\varepsilon_n$, so $y_0-\hat y(x_0)$ has mean zero and variance $\sigma^2+\Var\hat y(x_0)$.

\emph{The inequality.} For each $i$ the product $(U-x_i)(x_i-L)$ is nonnegative because $L\le x_i\le U$, and it is zero exactly when $x_i$ is $L$ or $U$. Summing over $i$ and expanding gives $\sum_ix_i^2\le(U+L)\sum_ix_i-nUL$; subtracting $n\bar x^2$ from both sides gives
\begin{equation}\label{eq:extrapolation:bd}
S_{xx}\le n(U-\bar x)(\bar x-L),
\end{equation}
with equality if and only if every $x_i$ equals $L$ or $U$.

\emph{Part (ii).} The product $(U-\bar x)(\bar x-L)$ equals $(U-L)^2/4-\bigl(\bar x-(L+U)/2\bigr)^2$, so it is at most $(U-L)^2/4$, with equality exactly when $\bar x$ is the midpoint. Equality throughout therefore holds exactly when all the design points sit at the two ends with their mean at the midpoint, that is, half at each end, which requires $n$ to be even. The bound on $\Var\hat b$ follows from part~(i).

\emph{Part (iii).} Write $h=U-L$ and $p=(\bar x-L)/h$. The design points are not all equal, so $S_{xx}>0$, and \eqref{eq:extrapolation:bd} forces $0<p<1$. Since $x_0-\bar x=(t-p)h$ and, by \eqref{eq:extrapolation:bd}, $S_{xx}\le np(1-p)h^2$, part~(i) gives
\[
\frac{n\Var\hat y(x_0)}{\sigma^2}=1+\frac{n(x_0-\bar x)^2}{S_{xx}}\ge1+\frac{(t-p)^2}{p(1-p)}.
\]
To compare the right side with $(2t-1)^2$, expand both sides of the identity $(t-p)^2-4t(t-1)\,p(1-p)=\bigl(t-(2t-1)p\bigr)^2$; since $(2t-1)^2-1=4t(t-1)$, it rearranges to
\[
1+\frac{(t-p)^2}{p(1-p)}=(2t-1)^2+\frac{\bigl(t-(2t-1)p\bigr)^2}{p(1-p)}\ge(2t-1)^2 .
\]
Equality requires equality in \eqref{eq:extrapolation:bd}, so that every design point is $L$ or $U$, and $p=t/(2t-1)$; for such a design $p$ is the fraction of points at $U$. Because $t>1$, this fraction lies strictly between $\tfrac12$ and $1$, so the design has points at both ends, and it attains the bound.
\end{proof}

\begin{corollary}[A prediction interval for one new run]\label{cor:extrapolation:interval}
In the setting of \cref{thm:extrapolation:variance}(i), suppose in addition that $n\ge3$ and that $\varepsilon_0,\varepsilon_1,\dots,\varepsilon_n$ are independent and normal, and let $S^2$ be the residual sum of squares of the fit divided by $n-2$. Then the interval
\[
\hat y(x_0)\pm t_{n-2,0.975}\,S\sqrt{1+\frac1n+\frac{(x_0-\bar x)^2}{S_{xx}}}
\]
contains $y_0$ with probability exactly $0.95$. Without the term $1$ under the root it is a 95\% confidence interval for the mean response $a+bx_0$.
\end{corollary}

\begin{proof}
For a linear model with independent normal errors, the residual sum of squares divided by $\sigma^2$ has the $\chi^2$ distribution with $n-2$ degrees of freedom and is independent of $(\hat a,\hat b)$ \parencite{SeberLee2003}. The error $\varepsilon_0$ is independent of both, so $y_0-\hat y(x_0)$ is normal, with mean zero and the variance of \cref{thm:extrapolation:variance}(i), and independent of $S^2$. A standard normal divided by the square root of an independent $\chi^2_{n-2}/(n-2)$ variable has Student's $t$ distribution with $n-2$ degrees of freedom, as in \cref{thm:toolkit:tinterval}, and the interval is this statement solved for $y_0$. The confidence interval follows in the same way without $\varepsilon_0$.
\end{proof}

\paragraph{Reading the theorem.} In logarithmic coordinates an interval of half-width $h$ for $\ln L$ is a factor $e^{\pm h}$ on the loss, and far from the ladder the half-width grows in proportion to the distance from its center. Distance is measured in units of the ladder's own spread: what matters is $(x_0-\bar x)^2/S_{xx}$, not the number of decades.\sidenote{In the language of regression diagnostics, $1/n+(x_0-\bar x)^2/S_{xx}$ is the leverage of the point $x_0$; a forecast beyond the ladder is a prediction at a leverage larger than that of its top rung.} Consider eight runs on a ladder spanning two decades and a forecast one decade beyond its top, so that $t=1.5$. Spread evenly, the runs give $n\Var\hat y(x_0)/\sigma^2=10.3$; four at each end give 5; and six at the top with two at the bottom, the fraction $t/(2t-1)=3/4$ of part~(iii), give the minimum, 4. The even ladder's forecast is thus 1.6 times as uncertain as the best one in standard deviation, and 3.2 times as uncertain as its own prediction at the center of the ladder. Two rules follow: the variance depends on the design only through $n$, $\bar x$ and $S_{xx}$, so runs in the middle of a ladder buy little precision; and a forecast upward is most precise when most runs sit at the top of the ladder.\sidenote{Part~(iii) is the simplest case of the optimal allocation of \textcite{Elfving1952}; \textcite{HoelLevine1964} found the optimal designs for extrapolating a polynomial of every degree.}

But a design with runs only at the two ends of the ladder cannot test the line it fits: two scales determine a line exactly, whatever the curve between and beyond them. The precision of \cref{thm:extrapolation:variance} is the precision of a forecast whose form is right, and the next section asks what happens when it is not.

\section{The method: the form, the break and the backtest}\label{sec:extrapolation:form}

A power law with an irreducible term, $L=E+AN^{-\alpha}$, is not a straight line in logarithmic coordinates: its local slope is $-\alpha(1-E/L)$, which flattens as the loss approaches its floor (\cref{ch:protocol}). The simplest model of a wrong form is a line fitted to a parabola.

\begin{proposition}[A line fitted to a curve]\label{prop:extrapolation:misspec}\index{misspecification}
Let the design be as in \cref{thm:extrapolation:variance} with $n\ge3$, and let $y_i=f(x_i)+\varepsilon_i$ with $f(x)=a+bx+cx^2$ and errors as there. Fit a straight line by least squares. Write $m_k=\frac1n\sum_i(x_i-\bar x)^k$ for $k=2,3,4$, and
\[
r(x)=(x-\bar x)^2-\frac{m_3}{m_2}(x-\bar x)-m_2,\qquad\kappa=m_4-m_2^2-\frac{m_3^2}{m_2}.
\]
\begin{enumerate}
\item For every $x_0$, $\E\hat y(x_0)=f(x_0)-c\,r(x_0)$, and $\Var\hat y(x_0)$ is given by \eqref{eq:extrapolation:var}.
\item The residuals have means $\E\bigl(y_i-\hat y(x_i)\bigr)=c\,r(x_i)$, and $\E S^2=\sigma^2+c^2n\kappa/(n-2)$.
\item If $c\ne0$, then for a new observation $y_0=f(x_0)+\varepsilon_0$, with $\varepsilon_0$ as in \cref{thm:extrapolation:variance}(i), the probability that the interval of \cref{cor:extrapolation:interval} contains $y_0$ tends to zero as $\abs{x_0}\to\infty$; the same holds for every fixed multiplier in place of $t_{n-2,0.975}$.
\end{enumerate}
\end{proposition}

\begin{proof}
The plan is to use that least squares is linear in the responses, so that the expected fit is the fit of $f$ itself, which we compute; part~(ii) then follows, and part~(iii) compares a bias that grows quadratically with an interval that grows linearly.

\emph{Part (i).} By the proof of \cref{thm:extrapolation:variance}, $\hat y(x_0)=\sum_i\ell_iy_i$ with weights $\ell_i=1/n+(x_0-\bar x)w_i$ that depend only on the design. Hence $\E\hat y(x_0)=\sum_i\ell_if(x_i)$, the value at $x_0$ of the least-squares line through the points $(x_i,f(x_i))$, and the random part of $\hat y(x_0)$ is the same combination of the errors as before, so its variance is \eqref{eq:extrapolation:var}. Fitting a line is linear in the responses and reproduces every linear function exactly. Write $u=x-\bar x$; then $cx^2=c\bar x^2+2c\bar xu+cu^2$, and only $cu^2$ is not reproduced. Since $\sum_iu_i=0$, the least-squares line through the points $(u_i,u_i^2)$ has intercept $\frac1n\sum_iu_i^2=m_2$ and slope $\sum_iu_i^3/\sum_iu_i^2=m_3/m_2$. The fit of $f$ at $x_0$ therefore falls short of $f(x_0)$ by $c\bigl((x_0-\bar x)^2-m_2-(m_3/m_2)(x_0-\bar x)\bigr)=c\,r(x_0)$.

\emph{Part (ii).} At a design point, $\E\bigl(y_i-\hat y(x_i)\bigr)=f(x_i)-\E\hat y(x_i)=c\,r(x_i)$ by part~(i). The residual vector is therefore the vector with entries $c\,r(x_i)$ plus the residuals that the errors alone would leave; the cross term has mean zero, and the residual sum of squares of the errors alone is $\varepsilon^\top(I-H)\varepsilon$, for the projection $I-H$ onto the orthogonal complement of the constant vector and the vector of design points. Because the errors are uncorrelated with variance $\sigma^2$, its mean is $\sigma^2$ times the trace of $I-H$, which is $n-2$. It remains to show $\sum_ir(x_i)^2=n\kappa$. Expanding the square of $r(x_i)=(u_i^2-m_2)-(m_3/m_2)u_i$ and using $\sum_i(u_i^2-m_2)^2=n(m_4-m_2^2)$, $\sum_iu_i(u_i^2-m_2)=nm_3$ and $\sum_iu_i^2=nm_2$,
\[
\sum_ir(x_i)^2=n(m_4-m_2^2)-2\frac{m_3}{m_2}nm_3+\frac{m_3^2}{m_2^2}nm_2=n\kappa .
\]
Dividing the mean of the residual sum of squares by $n-2$ gives the formula for $\E S^2$.

\emph{Part (iii).} Write $\delta=x_0-\bar x$ and $q=1+1/n+\delta^2/S_{xx}$, and let $Z=\bigl(\varepsilon_0-\hat y(x_0)+\E\hat y(x_0)\bigr)/(\sigma\sqrt q)$, which has mean zero and variance one by part~(i) and \cref{thm:extrapolation:variance}(i). By part~(i) the prediction error is $y_0-\hat y(x_0)=c\,r(x_0)+\sigma\sqrt q\,Z$, so for a multiplier $\tau>0$ the interval contains $y_0$ exactly when
\[
\bigl|c\,r(x_0)/\sqrt q+\sigma Z\bigr|\le\tau S .
\]
As $\abs{x_0}\to\infty$, $r(x_0)/\delta^2\to1$ and $\sqrt q/\abs\delta\to1/\sqrt{S_{xx}}$, so $A=\abs{c\,r(x_0)}/\sqrt q$ grows like $\abs c\abs\delta\sqrt{S_{xx}}$ and tends to infinity because $c\ne0$. The event above requires $\sigma\abs Z\ge A/2$ or $\tau S\ge A/2$. By Chebyshev's inequality the first has probability at most $4\sigma^2/A^2$. The second has a probability that tends to zero because $S$ is computed from the design points alone, so its distribution does not change with $x_0$, and it is finite. Both probabilities tend to zero.
\end{proof}

\paragraph{Invisible inside, large outside.} The function $r$ does two jobs. At the design points, $c\,r(x_i)$ is the pattern that the curvature leaves in the residuals, with mean square $c^2\kappa$ over the design, which adds $c^2n\kappa/(n-2)$ to the expected residual variance $S^2$; beyond the design, $c\,r(x_0)$ is the error of the forecast. The two scales are very different. For a design spread evenly over a window of width $h$, in the limit of many points, $m_2=h^2/12$, $m_3=0$ and $\kappa=h^4/180$, so the curvature shows in the residuals with root mean square $\abs c h^2/\sqrt{180}\approx0.075\abs c h^2$, while the forecast half a window past the edge of the design, at $\delta=h$, is off by $\abs c(h^2-h^2/12)\approx0.92\abs c h^2$, twelve times as much, and one window past the edge by 29 times as much. A curvature too small to detect in the residuals is large where the forecast is made. This is the case that \textcite{BoxDraper1959} analyzed in their classic paper on design: a line fitted where the truth is quadratic. Their term in the residual sum of squares that makes the curvature detectable is $c^2n\kappa$, and they concluded that when the line must describe the curve over a region, admitting the error of the form moves the best design from the ends of the region toward its middle, the opposite of the answer when the form is taken as right.

For scaling curves the sign is known. In logarithmic coordinates a law with a floor has second derivative $\alpha^2E(L-E)/L^2$, which is positive, so a line fitted to it predicts too low a loss beyond the ladder: the error of the form makes the forecast optimistic. The curvature is small among small models, whose loss lies far above its floor, and largest where $L=2E$; a ladder that approaches its floor is the one whose line misleads most. Fitting the right form helps only if it is the right form, and a different kind of error lies beyond every form.

\begin{remark}[A break beyond the data]\label{rem:extrapolation:break}
Data below a scale $U$ cannot rule out a change of slope above it. Take the smoothly broken power law of \textcite{CaballeroEtAl2023} with one break, $g(x)=bx^{-c_0}\bigl(1+(x/d)^{1/f}\bigr)^{-c_1f}$, where now $x>0$ is the scale itself, and let $U<d<x_0$. As the sharpness parameter $f$ tends to zero, the last factor tends to $1$ uniformly on $(0,U]$ and to $(x_0/d)^{-c_1}$ at $x_0$ (\cref{ex:extrapolation:break}). So for every $\eps>0$ there is a broken law that differs from the plain power law $bx^{-c_0}$ by less than a factor $1+\eps$ at every design point and, at $x_0$, by a factor within $1+\eps$ of $(x_0/d)^{-c_1}$, which is as large or as small as $c_1$ makes it. With normal errors the two laws give distributions of the data whose total variation distance tends to zero, so every test based on the data has power at most its level plus a term that tends to zero. A forecast at $x_0$ therefore rests on an assumption that the data cannot test: that no break lies between the ladder and the target.
\end{remark}

\paragraph{Backtesting.} The theorem and the proposition bound the two errors that can be computed. Whether a procedure's intervals are honest at a given distance is an empirical question about the curve, and it has an empirical answer.

\begin{definition}[Backtest and coverage]\label{def:extrapolation:backtest}\index{backtest}\index{coverage}
Let a ladder of rungs $x_1<\dots<x_m$ be measured. A \term{backtest} of a forecasting procedure fits the procedure to the rungs $x_1,\dots,x_k$ only and predicts each held-out rung $x_j$, $j>k$, with a nominal 95\% interval. The \term{coverage} at distance $x_j-x_k$ is the probability that the observed outcome at $x_j$ falls in its interval.
\end{definition}

Coverage is estimated by counting over independent replicates of the ladder, with an exact binomial interval (\cref{thm:toolkit:clopper}); backtests that share a ladder share its noise and are not independent. Three practices make a backtest informative. Fit several forms, at least a law with a floor and a form that allows a break, and report the spread of their forecasts as a lower bound on the uncertainty that the choice of form adds. Run a simulated control in which the assumed form holds exactly and the noise has the observed size: the positive control of \cref{ch:controls}, which separates the error of the interval's approximations, such as the delta method for a nonlinear fit (\cref{thm:toolkit:nls,thm:toolkit:delta}), from the error of the form. And state the \term{domain of validity}\index{domain of validity}: the family, the protocol and the prompts of the ladder, the distance of the target in decades, and the coverage measured at that distance. A backtest can only measure coverage at distances shorter than the ladder, so a forecast beyond the backtested distance carries the assumption of \cref{rem:extrapolation:break} in full. The forecasts that held in \cref{sec:extrapolation:postmortem} had these features: smooth quantities, a ladder trained like its target, and, for GPT-4, a forecast written down before the result was known.

\section{In the laboratory: a ladder of widths and a curve that turns}\label{sec:extrapolation:lab}

The first lab backtests forecasts on a smooth scaling curve, where parameter uncertainty and the error of the form are at work; the second fits short windows of a curve that turns, where the regime changes.

\begin{lab}{Backtesting a scaling curve}{13-extrap-backtest}\label{lab:extrapolation:backtest}
\labpart{Setup} The teacher task under the online protocol: one pass over 65,536 fresh examples, Adam at learning rate 0.01 (the best of 0.001, 0.003 and 0.01 at every width in a pilot) with a cosine schedule, batch 64. A ladder of 13 networks with two hidden layers of width 4 to 256 (61 to 68,353 parameters, 3.05 decades); 30 ladders, every run with its own seeds (390 runs). Each backtest fits the 5, 6, 7 or 8 smallest rungs of one ladder, by a straight line in $(\ln N,\ln L)$ with the interval of \cref{cor:extrapolation:interval} and by the law $\ln L=\ln(E+AN^{-\alpha})$ with a delta-method interval, and predicts every larger rung: 120 backtests and 780 forecasts per form. A simulated control repeats both on ladders drawn from a law with a floor. About three minutes.
\labpart{Predict first} Fitted to the six smallest rungs, by what factor will the line miss the largest? At what distance will its intervals stop covering? Will the law with a floor do better, and at what price in the width of its intervals?
\labpart{Look for} Where the curve bends toward its floor, relative to the fitted rungs; the coverage of each form against the distance; the width of the intervals; the curvature inside the window against the error beyond it.
\end{lab}

\begin{figure}[t]
\begin{fullwidth}
\centering
\begin{tikzpicture}
\begin{groupplot}[group style={group size=2 by 1,horizontal sep=17mm},lplot,
  width=0.5\linewidth,height=0.36\linewidth]
\nextgroupplot[title={(a) one backtest: the six smallest rungs fitted},xmode=log,ymode=log,lloglog,
  xlabel={parameters $N$},ylabel={test loss},xmin=45,xmax=95000,ymin=0.0015,ymax=0.1,
  xtick={100,1000,10000},xticklabels={$10^2$,$10^3$,$10^4$},
  ytick={0.002,0.005,0.01,0.02,0.05,0.1},yticklabels={$0.002$,$0.005$,$0.01$,$0.02$,$0.05$,$0.1$},
  legend style={at={(0.03,0.03)},anchor=south west,font=\scriptsize},legend cell align=left]
\addplot[name path=sh,draw=none,forget plot] table[x=N,y=sathi] {figures/data/13-extrapolation-backtest-fit.dat};
\addplot[name path=sl,draw=none,forget plot] table[x=N,y=satlo] {figures/data/13-extrapolation-backtest-fit.dat};
\addplot[fill=inkblue!11,draw=none,forget plot] fill between[of=sh and sl];
\addplot[name path=lh,draw=none,forget plot] table[x=N,y=linehi] {figures/data/13-extrapolation-backtest-fit.dat};
\addplot[name path=ll,draw=none,forget plot] table[x=N,y=linelo] {figures/data/13-extrapolation-backtest-fit.dat};
\addplot[fill=rubric!16,draw=none,forget plot] fill between[of=lh and ll];
\addplot[reference,forget plot] coordinates {(45,0.01) (95000,0.01)};
\addplot[run,forget plot] table[x=N,y=loss] {figures/data/13-extrapolation-backtest-runs.dat};
\addplot[mean,ebars,forget plot] table[x=N,y=mean,y error=hw] {figures/data/13-extrapolation-backtest-means.dat};
\addplot[fit,unbounded coords=jump] table[x=N,y=linein] {figures/data/13-extrapolation-backtest-fit.dat};
\addlegendentry{straight line}
\addplot[extrapolation,unbounded coords=jump,forget plot] table[x=N,y=lineout] {figures/data/13-extrapolation-backtest-fit.dat};
\addplot[inkblue,line width=1.0pt,no marks,unbounded coords=jump] table[x=N,y=satin] {figures/data/13-extrapolation-backtest-fit.dat};
\addlegendentry{law with a floor}
\addplot[inkblue,line width=0.9pt,dash pattern=on 3pt off 2pt,no marks,unbounded coords=jump,forget plot] table[x=N,y=satout] {figures/data/13-extrapolation-backtest-fit.dat};
\addplot[only marks,mark=diamond*,mark size=2.3pt,color=ink,mark options={fill=ink,draw=white,line width=0.2pt},unbounded coords=jump,forget plot] table[x=N,y=fitted] {figures/data/13-extrapolation-backtest-ladder0.dat};
\addplot[only marks,mark=diamond*,mark size=2.3pt,color=ink,mark options={fill=white,line width=0.5pt},unbounded coords=jump,forget plot] table[x=N,y=heldout] {figures/data/13-extrapolation-backtest-ladder0.dat};
\nextgroupplot[title={(b) coverage of the 95\% intervals, 30 ladders},xlabel={decades beyond the largest fitted rung},
  ylabel={coverage},xmin=0,xmax=2.08,ymin=0,ymax=1.06,xtick={0,0.5,1,1.5,2},
  ytick={0,0.25,0.5,0.75,0.95},yticklabels={$0$,$0.25$,$0.5$,$0.75$,$0.95$},
  legend style={at={(0.03,0.03)},anchor=south west,font=\scriptsize},legend cell align=left]
\addplot[reference,forget plot] coordinates {(0,0.95) (2.08,0.95)};
\addplot[meanB,error bars/.cd,y dir=both,y explicit,error bar style={rubric,line width=0.5pt},
  error mark options={rubric,rotate=90,mark size=1.7pt,line width=0.5pt}]
  table[x=dist,y=covl,y error plus expr=\thisrow{hil}-\thisrow{covl},y error minus expr=\thisrow{covl}-\thisrow{lol}] {figures/data/13-extrapolation-backtest-coverage-k6.dat};
\addlegendentry{straight line}
\addplot[mean,ebars] table[x expr=\thisrow{dist}+0.035,y=covs,y error plus expr=\thisrow{his}-\thisrow{covs},y error minus expr=\thisrow{covs}-\thisrow{los}] {figures/data/13-extrapolation-backtest-coverage-k6.dat};
\addlegendentry{law with a floor}
\addplot[rubric,line width=0.7pt,densely dotted,no marks] table[x=dist,y=ctll] {figures/data/13-extrapolation-backtest-coverage-k6.dat};
\addlegendentry{simulated control}
\addplot[inkblue,line width=0.7pt,densely dotted,no marks,forget plot] table[x=dist,y=ctls] {figures/data/13-extrapolation-backtest-coverage-k6.dat};
\end{groupplot}
\end{tikzpicture}
\caption{\Cref{lab:extrapolation:backtest}: the teacher task, one pass over 65,536 examples, two hidden layers of 13 widths from 4 to 256, 30 ladders with seeds $1000r+j$ for ladder $r$ and rung $j$. (a)~Hollow marks: all 390 runs; filled circles: means over the ladders with 95\% $t$-intervals; dashed grey line: the noise floor of the teacher, 0.01. The two forms are fitted to the six smallest rungs of ladder~0 (filled diamonds; hollow diamonds are its held-out rungs) and drawn solid over the fitted rungs and dashed beyond, with bands for the 95\% prediction interval of one run: vermilion for the straight line, blue for the law with a floor. (b)~The same window in all 30 ladders: the fraction of held-out runs inside their intervals, with 95\% Clopper--Pearson intervals over ladders (offset for legibility); dotted lines: the same procedures on 600 simulated ladders in which a law with a floor holds exactly.}
\label{fig:extrapolation:backtest}
\end{fullwidth}
\end{figure}

\paragraph{What happened.} The mean test loss falls from \pmval{0.0397}{0.0017} at 61 parameters to \pmval{0.0122}{0.0001} at 34,753 and \pmval{0.0123}{0.0001} at 68,353 (95\% $t$-intervals over the 30 ladders), above the teacher's noise floor of 0.01 by the part that 65,536 examples cannot remove.\sidenote{In the notation of \cref{ch:fitting}, a curve in $N$ at fixed $D$ has the floor $E+A_DD^{-\alpha_D}$, not $E$, and the law with a floor fits that sum.} In logarithmic coordinates the local slope is $-0.33$ to $-0.37$ over the first four rungs and between $-0.04$ and $+0.01$ over the last four: the curve flattens toward its floor between about $10^3$ and $10^4$ parameters, beyond the six smallest rungs. The spread of $\ln L$ over runs falls from 0.11 at the smallest width to 0.03 at the largest.

In ladder~0 (\cref{fig:extrapolation:backtest}a), the line fitted to the six smallest rungs predicts a loss of 0.0041 at 68,353 parameters, with a 95\% interval of 0.0023 to 0.0074 that lies wholly below the teacher's noise floor; the run there scored 0.0125. The law with a floor predicts 0.0110 with an interval of 0.0014 to 0.085, a factor of 60 wide: honest, and of little use. Across ladders its fitted floor is barely identified, with median 0.0107 and quartiles 0.0001 and 0.0139 (\cref{ch:fitting}).

Over all ladders and windows, the line's intervals cover 91\% of the held-out runs within 0.55 decades of the window (87--96\%, a 95\% bootstrap interval over whole ladders), 41\% (34--48\%) between 0.55 and 1.2 decades, and 21\% (13--29\%) beyond; the law with a floor covers 99\%, 93\% and 88\% (82--94\%). The simulated control reproduces both patterns (\cref{fig:extrapolation:backtest}b): for the window of six rungs the line's coverage falls from 94\% to 7\% there too, so its failure is a failure of the form, and the law with a floor falls from 96\% to 78\%, further than on the real ladders, so its shortfall at long range comes from its delta-method interval, linear in parameters that six rungs barely determine, and not from a departure of the teacher's curve from the form.

\begin{table}[t]
\centering\small\setlength{\tabcolsep}{3pt}
\begin{tabular}{l*{7}{r}}
\toprule
decades beyond the window & 0.25 & 0.51 & 0.79 & 1.07 & 1.36 & 1.65 & 1.94\\
\midrule
\emph{straight line}\\
\quad error, observed & $-0.08$ & $-0.16$ & $-0.29$ & $-0.44$ & $-0.63$ & $-0.81$ & $-1.03$\\
\quad its 95\% half-width & $0.03$ & $0.03$ & $0.04$ & $0.04$ & $0.05$ & $0.06$ & $0.07$\\
\quad predicted, $-c\,r(x_0)$ & $-0.07$ & $-0.15$ & $-0.25$ & $-0.38$ & $-0.53$ & $-0.71$ & $-0.91$\\
\quad interval half-width & $0.28$ & $0.32$ & $0.37$ & $0.42$ & $0.47$ & $0.53$ & $0.59$\\
\emph{law with a floor}\\
\quad error, observed & $+0.02$ & $+0.02$ & $-0.02$ & $-0.06$ & $-0.13$ & $-0.19$ & $-0.27$\\
\quad interval half-width & $0.39$ & $0.51$ & $0.60$ & $0.68$ & $0.77$ & $0.86$ & $0.94$\\
\bottomrule
\end{tabular}
\caption{\Cref{lab:extrapolation:backtest}, window of the six smallest rungs: errors of the forecast of $\ln L$ (forecast minus observed, in natural logarithms) averaged over the 30 ladders, with the 95\% half-width of that average, and median half-widths of the 95\% prediction intervals of one run. The row $-c\,r(x_0)$ is the error that \cref{prop:extrapolation:misspec} predicts from the curvature $c=0.028$ of a parabola fitted to the mean of $\ln L$ over the six rungs.}
\label{tab:extrapolation:bias}
\end{table}

\Cref{tab:extrapolation:bias} tests \cref{prop:extrapolation:misspec}. A parabola fitted to the mean of $\ln L$ over the six smallest rungs has curvature $c=0.028$ per unit of $(\ln N)^2$ (95\% interval over ladders 0.009--0.048), and with it the proposition predicts an error of $-0.38$ at 1.07 decades and $-0.91$ at 1.94, against \pmval{-0.44}{0.04} and \pmval{-1.03}{0.07} observed. The point estimate accounts for 84--97\% of the error at the seven distances, and over the interval of the curvature the prediction brackets the observed error at every distance, here and for the windows of seven and eight rungs. The line's half-width grows from 0.28 to 0.59 while its error grows from 0.08 to 1.03, and the error passes the half-width near one decade, where the coverage falls through one half. Inside the window the same curvature is invisible. It adds $c^2n\kappa/(n-2)=0.0005$ to the residual variance, against 0.0056 for the spread of runs; from a single ladder its standard error is 0.047, larger than the curvature itself; and a $t$-test of a quadratic term rejects the line at level 0.05 in 3 of the 30 ladders. For the window of five rungs the curvature is not determined even by all 30 ladders (interval $-0.023$ to $0.047$), and the proposition predicts nothing useful.

\begin{lab}{A curve that turns}{13-extrap-ushape}\label{lab:extrapolation:ushape}
\labpart{Setup} Digits with 20\% of the training labels replaced by random labels; a network with one hidden layer of width 2 to 2048 (160 to 153,610 parameters), Adam at learning rate $10^{-3}$, batch 64, 200 epochs, no regularization; six runs per width, run $s$ with seed $s$ for initialization, order, split and label noise (108 runs). Two windows are fitted by straight lines in $(\ln N,\ln\text{error})$ and extrapolated with the intervals of \cref{cor:extrapolation:interval}: window A on widths 4--12, window B on widths 16--64. About three minutes.
\labpart{Predict first} The training set has 1,258 images, about 226 of them with wrong labels. At what width will the networks fit every training label? What will the test error do as the width grows past that point?
\labpart{Look for} The two descents and the peak between them; the training error at the peak; the local slope over short windows, and how its sign depends on where the window sits.
\end{lab}

\begin{figure}[t]
\begin{fullwidth}
\centering
\begin{tikzpicture}
\begin{groupplot}[group style={group size=2 by 1,horizontal sep=17mm},lplot,
  width=0.5\linewidth,height=0.34\linewidth,xmode=log,xmin=120,xmax=230000,
  xtick={1000,10000,100000},xticklabels={$10^3$,$10^4$,$10^5$},xlabel={parameters $N$}]
\nextgroupplot[title={(a) test error, and two windows extrapolated},ylabel={error (\%)},ymin=0,ymax=56,
  ytick={0,10,20,30,40,50},legend style={at={(0.22,0.97)},anchor=north west,font=\scriptsize},legend cell align=left]
\addplot[name path=ah,draw=none,forget plot] table[x=N,y=hi] {figures/data/13-extrapolation-ushape-fitA.dat};
\addplot[name path=al,draw=none,forget plot] table[x=N,y=lo] {figures/data/13-extrapolation-ushape-fitA.dat};
\addplot[fill=inkblue!11,draw=none,forget plot] fill between[of=ah and al];
\addplot[name path=bh,draw=none,forget plot] table[x=N,y=hi] {figures/data/13-extrapolation-ushape-fitB.dat};
\addplot[name path=bl,draw=none,forget plot] table[x=N,y=lo] {figures/data/13-extrapolation-ushape-fitB.dat};
\addplot[fill=rubric!16,draw=none,forget plot] fill between[of=bh and bl];
\addplot[run,forget plot] table[x=N,y=err] {figures/data/13-extrapolation-ushape-runs.dat};
\addplot[mean,ebars] table[x=N,y=mean,y error=hw] {figures/data/13-extrapolation-ushape-means.dat};
\addlegendentry{test error}
\addplot[black!55,line width=0.6pt,mark=square*,mark size=0.9pt,mark options={solid}] table[x=N,y=train] {figures/data/13-extrapolation-ushape-means.dat};
\addlegendentry{training error}
\addplot[inkblue,line width=1.0pt,no marks,unbounded coords=jump,forget plot] table[x=N,y=predin] {figures/data/13-extrapolation-ushape-fitA.dat};
\addplot[inkblue,line width=0.9pt,dash pattern=on 3pt off 2pt,no marks,unbounded coords=jump] table[x=N,y=predout] {figures/data/13-extrapolation-ushape-fitA.dat};
\addlegendentry{window A}
\addplot[fit,unbounded coords=jump,forget plot] table[x=N,y=predin] {figures/data/13-extrapolation-ushape-fitB.dat};
\addplot[extrapolation,unbounded coords=jump] table[x=N,y=predout] {figures/data/13-extrapolation-ushape-fitB.dat};
\addlegendentry{window B}
\nextgroupplot[title={(b) local slope over three widths},ylabel={$\dd\ln(\text{error})/\dd\ln N$},ymin=-1.9,ymax=0.8,
  ytick={-1.5,-1,-0.5,0,0.5}]
\addplot[reference] coordinates {(120,0) (230000,0)};
\addplot[mean,ebars] table[x=N,y=slope,y error plus expr=\thisrow{hi}-\thisrow{slope},y error minus expr=\thisrow{slope}-\thisrow{lo}] {figures/data/13-extrapolation-ushape-slopes.dat};
\end{groupplot}
\end{tikzpicture}
\caption{\Cref{lab:extrapolation:ushape}: digits with 20\% of the training labels random, one hidden layer of widths 2 to 2048, six runs per width (seeds 0--5). (a)~Hollow marks: single runs; filled circles: means with 95\% $t$-intervals; small squares: mean training error on the noisy training labels. Straight lines in $(\ln N,\ln\text{error})$ fitted to all runs in window A (widths 4--12, blue) and window B (widths 16--64, vermilion), solid over the window and dashed beyond, with bands for the 95\% prediction interval of one run (window B's band is cut by the frame). (b)~The slope of a line fitted to the runs of three consecutive widths, plotted at their geometric mean, with 95\% intervals.}
\label{fig:extrapolation:ushape}
\end{fullwidth}
\end{figure}

\paragraph{What happened.} The mean test error falls from \pmval{50.2}{7.5}\% at width 2 to \pmval{10.4}{2.1}\% at width 8 and \pmval{10.5}{1.2}\% at width 12, rises to \pmval{17.1}{1.6}\% at width 64, where the networks still miss 1.8\% of their training labels, and falls again once they fit nearly all of them, from width 96 on: to \pmval{13.0}{1.4}\% at width 512 and \pmval{12.5}{0.9}\% at 2048, still above the first minimum, as the largest PaLM model was still below the smallest on four of the six U-shaped tasks of \cref{fd:extrapolation:ushape} (\cref{fig:extrapolation:ushape}a). This is the \term{double descent}\index{double descent} of \textcite{BelkinEtAl2019} and \textcite{NakkiranEtAl2020}, with its peak where the network can just fit the noisy training set.\sidenote{The second volume of this series, \emph{The Hydrogen Atoms of Deep Learning}, derives the double-descent curve of least squares exactly in its first chapter.}

Window B is a miniature inverse-scaling trend: over widths 16 to 64 the error rises with slope 0.29 per unit of $\ln N$ (95\% interval 0.19--0.38), a factor of 1.9 per decade, and no test inside the window would doubt it. Extrapolated, it predicts 20.7\% at width 128, with a 95\% interval for one run of 15.5--27.6\%, and 45.7\% (28.7--72.9\%) at width 2048; the runs scored 15.0\% and 12.5\% on average, and of the 42 runs beyond the window only 7 fall inside their intervals, all at the two widths next to it. Window A, with slope $-0.56$ ($-0.79$ to $-0.34$), extrapolates the first descent to 3.6\% at width 64 and 0.5\% at width 2048; 5 of the 72 runs beyond it are covered, all at the next width. A law with a floor fitted to window A does no better: it predicts 10.3\% at every larger width, below the peak and below the second descent. The local slope (\cref{fig:extrapolation:ushape}b) is $-1.41$ (95\% interval $-1.74$ to $-1.07$) near 230 parameters, positive between about 1,250 and 3,500 parameters ($+0.41$, interval $+0.18$ to $+0.64$, near 1,740), negative again near 7,000 ($-0.19$, interval $-0.37$ to $-0.01$), and within 0.11 of zero beyond 13,000. A short window estimates a local slope, which is all that a short window can estimate, and its sign depends on where the window sits.

The two labs reproduce the three errors of a forecast: the line's intervals widen with distance as \cref{thm:extrapolation:variance} says, its error from curvature grows as \cref{prop:extrapolation:misspec} says and escapes its intervals within a decade, and a change of regime defeats every form fitted before it. The second miniature reproduces the shape of the post-mortem, not its mechanism: its peak comes from fitting noisy labels, not from a distractor task.

\section{Protocol}\label{sec:extrapolation:protocol}

\begin{protocol}{Before you extrapolate}\label{prot:extrapolation:before}
\item State the target: the scale, the quantity and the protocol of the run to be forecast; write the forecast and its interval down, with the date, before the run finishes.
\item Design the ladder: span as many decades as the budget allows, put most runs at the top and the bottom of the range and a few in between to test the form, and repeat some rungs with new seeds to measure the noise apart from the curvature.
\item Train every rung by the protocol of the target: the same data, schedule matched to each horizon, hyperparameters scaled by a fixed rule (\cref{ch:protocol}).
\item Forecast a smooth quantity, such as a loss, the log-likelihood of the correct answer or an average over tasks, and reach a thresholded metric through it (\cref{ch:emergence}).
\item Fit at least two forms, one with a floor and one that allows a break, with intervals for one new run that include the uncertainty of the parameters (\cref{thm:extrapolation:variance}; \cref{thm:toolkit:delta,thm:toolkit:bootstrap}); report the spread between forms.
\item Backtest each form on held-out upper rungs over independent ladders, and report coverage and error against distance beside a simulated control in which the form holds.
\item State the domain of validity: family, protocol, prompts, range of the ladder, the distance of the target, and the coverage measured at that distance; say that a target beyond the backtested distance assumes that no break intervenes.
\item After the run, publish the comparison of forecast and result, whichever way it came out.
\end{protocol}

The protocol bounds the two errors that can be computed and checks the error of the form over the distances a backtest can reach. It does not protect against a change of regime beyond the ladder, which no data below the change can reveal (\cref{rem:extrapolation:break}); a backtest at one decade says nothing about three; and it does not make a forecast of one task's thresholded metric reliable, because a small change of prompt or harness can change the trend itself (\cref{sec:extrapolation:postmortem}). It is also expensive: replicate ladders multiply the cost of a ladder, whose cost is dominated by its top rungs. When the budget allows one run per rung, report the residual spread of the fit, which mixes the noise with the curvature (\cref{prop:extrapolation:misspec}), as an upper bound on the noise, give the forecast as a range across forms, and prefer decisions that need only a ranking at small scale: in a study of 25 pretraining corpora, the ranking of models at a single small size picked the better corpus at a scale of one billion parameters in about 80\% of comparisons, and none of eight scaling-law methods did better for the same compute \parencite{MagnussonEtAl2025}.

\begin{exercises}
\item \consolidation (Intervals on the loss.) (a) Show that an interval $\hat y\pm h$ for $\ln L$ is the interval $[\hat Le^{-h},\hat Le^{h}]$ for $L$, with $\hat L=e^{\hat y}$, and find the factor that the half-width 0.59 of \cref{tab:extrapolation:bias} allows. (b) Suppose that $\ln L$ is normal with mean $\mu$ and variance $s^2$. Show that $e^{\mu}$ is the median of $L$ and $e^{\mu+s^2/2}$ its mean, and say which one $e^{\hat y}$ estimates.\label{ex:extrapolation:loglog}

\item \consolidation (Three scales.) Three runs can be placed on a ladder spanning $[L,U]$ in logarithmic scale, to forecast at $x_0=L+2(U-L)$. (a) Show that two runs at $U$ and one at $L$ attain the bound of \cref{thm:extrapolation:variance}(iii), and compare with the design $L,(L+U)/2,U$. (b) A run at $U$ costs ten runs at $L$, and the budget is 31 runs at $L$. Which allocation between the two ends minimizes $\Var\hat y(x_0)$, and why does it differ from part~(a)?\label{ex:extrapolation:three}

\item \consolidation (Residuals and bias.) (a) Show that $\kappa\ge0$ in \cref{prop:extrapolation:misspec}, with equality if and only if the design has at most two distinct points, and conclude that a two-point design cannot reveal curvature. (b) Show that $\kappa=h^4/180$ for a design spread evenly over a window of width $h$ (in the limit of many points), and check the ratios 12 and 29 of \cref{sec:extrapolation:form}. (c) Show that $r$ has two real roots, and interpret them.\label{ex:extrapolation:residual}

\item \consolidation\hard (Breaks are invisible.) For $g_f(x)=bx^{-c_0}\bigl(1+(x/d)^{1/f}\bigr)^{-c_1f}$ with $0<U<d<x_0$, prove that $\sup_{0<x\le U}\abs{\ln g_f(x)-\ln(bx^{-c_0})}\to0$ and $\ln g_f(x_0)-\ln(bx_0^{-c_0})\to-c_1\ln(x_0/d)$ as $f\to0$. Then show that for normal errors of variance $\sigma^2$ in $\ln L$ at $n$ design points, the total variation distance between the distributions of the data under the two laws is at most $\norm{\mu-\mu'}/(\sigma\sqrt{2\pi})$, where $\mu$ and $\mu'$ are the two vectors of means of $\ln L$. \hint{use $\ln(1+s)\le s$ for the first limit, and for the bound on total variation reduce to one dimension along $\mu-\mu'$.}\label{ex:extrapolation:break}

\item \labwork (The ends of the ladder.) In \script{13-extrap-backtest.py}, replace the window of six rungs by a design with three runs at each of its two end rungs, drawn from three ladders. Predict from \cref{thm:extrapolation:variance}(ii)--(iii) how the width of the line's interval changes, then measure its coverage. Why can the law with a floor not be fitted to this design?\label{ex:extrapolation:ends}

\item \labwork (Weights and resampling.) The spread of $\ln L$ over runs falls from 0.11 to 0.03 along the ladder. Refit the line by weighted least squares with weights equal to the inverse variance of each rung, and replace the delta-method interval of the law with a floor by a parametric bootstrap (\cref{thm:toolkit:bootstrap}). Observe the coverage of both in the backtest and in the simulated control, and explain which changes.\label{ex:extrapolation:weights}

\item \labwork (Less noise, more noise.) Run \script{13-extrap-ushape.py} with 10\% and with 40\% of the training labels random. Predict where the peak moves and how high it rises before running, then check whether window B, refitted between the first minimum and the peak, still predicts an error that rises without end.\label{ex:extrapolation:noise}

\item \design (A registered forecast.) A team will train a model with $10^{24}$ FLOPs and wants its validation loss and its accuracy on one multiple-choice benchmark, with 1\% of that compute for the ladder. Design the ladder (sizes, tokens, seeds), the forms, the backtest and the registration document: what is written down before the run, and what result would count as a failed forecast.\label{ex:extrapolation:register}

\item \design (Testing for inverse scaling.) You have one family of six sizes and a task of 500 items. Design an evaluation that decides whether the task scales inversely over the range measured: prompts, items, checkpoints or seeds, the estimate of the slope and its interval, and the sentence that states what the result does and does not say about larger models.\label{ex:extrapolation:inverse}

\item \critique (Two decades past the data.) A plot shows six model sizes spanning one decade, a line fitted in log--log coordinates, and the line continued to a model a hundred times larger, labelled ``projected'', without an interval. Using \cref{thm:extrapolation:variance}, compute a lower bound on the standard deviation of the forecast relative to the noise of one run, and list what the plot would need to be an honest forecast.\label{ex:extrapolation:critique1}

\item \critique (Scaling hurts.) A paper evaluates four models of one family, with 1, 3, 7 and 13 billion parameters, on 200 items with one prompt; the accuracies are 61, 58, 57 and 54\%, and the paper concludes that scaling hurts the ability measured. Audit the claim: the binomial noise of 200 items, one model per size, one prompt, and the step to larger models.\label{ex:extrapolation:critique2}

\item \research (Detecting a break before it.) Tests of structural change find a break inside the data \parencite{Andrews1993}; a forecast needs to know whether a break lies beyond it. Formulate the problem of bounding, from a ladder below $U$, the probability of a break between $U$ and $x_0$ under a prior over breaks, and propose observable signals that might precede a break in deep learning, such as a drifting local slope, a growing spread over seeds, or the behavior of a subtask that acts as a distractor. Open: whether such signals predict the U-shaped tasks of \cref{fd:extrapolation:ushape} from the smaller models alone.\label{ex:extrapolation:research}
\end{exercises}

\notesandrefs
Scaling curves were forecasts from the start. \Textcite{HestnessEtAl2017} drew learning curves in four domains with three regions, small-data, power-law and irreducible-error, so that the power law described only the middle. \Textcite{KaplanEtAl2020} saw no deviation from their power laws, and noted that two of them, extrapolated, would contradict each other near $10^4$ petaflop-days, so that at least one had to break before that point. \Textcite{HoffmannEtAl2022} tested their forecast of the compute-optimal allocation by training Chinchilla, and the GPT-4 report turned such a forecast into a registered prediction \parencite{OpenAI2023}.

The statistics of forecasting with a fitted line is classical: prediction intervals in regression are in every text \parencite{SeberLee2003}, the design that minimizes the variance of an extrapolated prediction goes back to \textcite{Elfving1952} and, for polynomials of every degree, to \textcite{HoelLevine1964}, and \textcite{BoxDraper1959} showed that admitting an error of the form changes the best design. Econometrics treats changes of regime inside a series as structural breaks, with tests for a break at an unknown point \parencite{Andrews1993}.

In deep learning, \textcite{AlabdulmohsinNeyshaburZhai2022} argued that scaling laws should be judged by their error on held-out larger scales rather than by their fit, and released a benchmark of learning curves for doing so; \textcite{CaballeroEtAl2023} used it for the broken power law. For downstream tasks, \textcite{LourieHuCho2025} documented how often trends are irregular, while \textcite{GadreEtAl2025}, \textcite{BhagiaEtAl2025} and \textcite{RuanMaddisonHashimoto2024} showed how forecasts of averages and of losses can succeed; \textcite{MagnussonEtAl2025} found that for choosing pretraining data a comparison at one small scale is hard to beat. Double descent, the non-monotone curve of \cref{lab:extrapolation:ushape}, was named by \textcite{BelkinEtAl2019} and studied in deep networks, in model size, training time and data, by \textcite{NakkiranEtAl2020}.

\paragraph{Translating the literature.} McKenzie et al.'s inverse scaling is a negative local slope of a task's performance against the logarithm of training compute over the range measured; Wei et al.'s U-shaped and McKenzie et al.'s inverted-U scaling are curves whose local slope changes sign inside it. GPT-4's $L(C)=aC^b+c$ is our $E+AC^{-\gamma}$ with $c=E$, $a=A$ and $b=-\gamma$. In the broken power law of \textcite{CaballeroEtAl2023}, $y=a+bx^{-c_0}\prod_i(1+(x/d_i)^{1/f_i})^{-c_if_i}$, the constants $a$, $b$ and $c_0$ are our $E$, $A_N$ and $\alpha_N$ when $x$ is $N$; $c_i$ is the change of slope at the $i$th break, $d_i$ its location and $f_i$ its sharpness parameter (smaller is sharper), and \cref{rem:extrapolation:break} uses one break with $f\to0$. What \textcite{AlabdulmohsinNeyshaburZhai2022} call the extrapolation loss is the error of a backtest; a ladder of \textcite{BhagiaEtAl2025} is our design; the token multiplier $M$ of \textcite{GadreEtAl2025} is $D/N$; and the irreducible-error region of \textcite{HestnessEtAl2017} is the floor $E$.

\printchapterbib
\end{refsection}
